% Adapted from "Plantilla para Cartas" by Fausto Mauricio Lagos Suarez.
% Source: https://www.overleaf.com/latex/templates/plantilla-para-cartas/myxtzxtsmsdy
% Original template copyright 2006--2013 Xavier Danaux.
% Original source distributed under the LaTeX Project Public License v1.3c.
% Overleaf lists the template under Creative Commons Attribution 4.0.
% Changes: English language, personal and recipient details, no photo,
% and an original cover-letter draft for the MOTOR Ai application.
% Select cover_letter.tex as the Main document in Overleaf to compile.

\documentclass[11pt,a4paper,sans]{moderncv}
\moderncvstyle{casual}
\moderncvcolor{blue}
\nopagenumbers{}

\usepackage[T1]{fontenc}
\usepackage[utf8]{inputenc}
\usepackage[english]{babel}
\usepackage[scale=0.75]{geometry}
\usepackage{ragged2e}
\usepackage{microtype}

\name{Jin}{Tao}
\address{Elsterstr.\ 1}{14513 Teltow}{Germany}
\email{j.tao@campus.tu-berlin.de}
\homepage{github.com/JinTTTT}

\begin{document}

\recipient{MOTOR Ai Hiring Team}{Scheringstra\ss e 1\\13355 Berlin, Germany}
\date{\today}
\opening{Dear Hiring Team,}
\closing{Kind regards,}
\makelettertitle

\textbf{Application for AI/ML Engineer (m/f/d) Autonomous Systems}

\justifying
\setlength{\parindent}{0pt}

I am applying for the AI/ML Engineer position in Autonomous Systems at MOTOR Ai. I hold a master's degree in Vehicle Engineering from TU Berlin and currently work full-time at ZEKI. I want to bring my experience in model development, evaluation, and robotics integration to autonomous driving, contributing to your work toward deploying Level 4 systems.

In a visuo-tactile sensor-fusion project, I developed training and evaluation infrastructure for estimating in-hand object pose from RGB images and tactile measurements. I implemented and compared ten existing methods, then developed a transformer-based fusion model that matched the best evaluated baseline's pose-estimation accuracy. This work gave me experience translating published methods into working implementations and comparing models systematically. My technical toolkit includes Python, PyTorch, and TensorFlow.

At ZEKI, I developed a TurtleBot4 platform integrating SLAM, navigation, semantic mapping, and vision-language models. I also developed a Qwen-based agent that connects natural-language requests to navigation, exploration, and visual question answering. I am now developing a vision-language-action navigation policy integrated with Nav2, strengthening my experience connecting AI components to robot execution.

Alongside this work, I independently develop \textbf{TaoRobot: Autonomous Driving Full Stack}, an ongoing open-source personal project started in June 2026. I implement pose-graph SLAM, particle-filter localization, A* planning, and pure pursuit from scratch in lightweight ROS 2 nodes and integrate them in Gazebo. This gives me practical understanding of the autonomy stack into which learned models must fit. The code and demo are available at \href{https://github.com/JinTTTT/taorobot}{github.com/JinTTTT/taorobot}.

My deployment experience also includes YOLO-based gesture detection on a real autonomous shuttle and GPU-accelerated visual SLAM on NVIDIA Jetson Xavier. These projects connect my AI work with vehicle behavior and constrained onboard computing.

I am available to start on December 1, 2026. Only the thesis remains in my second master's degree in Computational Engineering Science; I complete it outside working hours, and it will not affect my full-time availability. I would welcome the opportunity to discuss my contribution to your team.

\makeletterclosing

\end{document}
